% openout_any=1 bibtex .latex/diploma

\documentclass[12pt,a4paper,openany]{book}

\usepackage[utf8]{inputenc}
\usepackage{cmap}
\usepackage{type1ec}
\usepackage[T1]{fontenc}
\usepackage{fancyhdr}
\usepackage{graphicx,epsfig}
\usepackage[slovene]{babel}
\usepackage{cite}
\usepackage{mathrsfs}
\usepackage{enumitem}
%\usepackage[nolist,nohyperlinks]{acronym}
\usepackage[acronym,toc,xindy]{glossaries}
\usepackage{todonotes}
\usepackage{mdframed}
\usepackage{mathtools}
\usepackage{amsmath}
\usepackage{amssymb}
\renewcommand{\baselinestretch}{1.3}

\usepackage[pdftex,colorlinks,citecolor=black,filecolor=black,linkcolor=black,urlcolor=black,pagebackref]{hyperref}
\usepackage{tikz}

%Velikost strani - dvostransko
\oddsidemargin 1.4cm
\evensidemargin 0.35cm
\textwidth 14cm
\topmargin 0.26cm
\headheight 0.6cm
\headsep 1.5cm
\textheight 20cm

%Nastavitev glave in repa strani
\pagestyle{fancy}
\fancyhead{}
\renewcommand{\chaptermark}[1]{\markboth{\textsf{Poglavje \thechapter:\ #1}}{}}
\renewcommand{\sectionmark}[1]{\markright{\textsf{\thesection\  #1}}{}}
\fancyhead[RE]{\leftmark}
\fancyhead[LO]{\rightmark}
\fancyhead[LE,RO]{\thepage}
\fancyfoot{}
\renewcommand{\headrulewidth}{0.0pt}
\renewcommand{\footrulewidth}{0.0pt}

\newcommand{\gnuplot}{\textbf{gnuplot}}
\newcommand{\pgfname}{\textsc{pgf}}
\newcommand{\tikzname}{Ti\emph{k}Z}

\DeclareMathAlphabet{\mathpzc}{OT1}{pzc}{m}{it}
\mdfsetup{nobreak=true}
\makeglossaries
%%%%%%%%%%%%%%%%%%%%%%%%%%%%%%%%%%%%%%%%%%%%%%%%%%%%%%%%%%%%%%%%
%% ccBeamer 0.1, 2007-07-02                                   %%
%% Written by Sebastian Pipping <webmaster@hartwork.org>      %%
%% ---------------------------------------------------------- %%
%% Licensed under Creative Commons Attribution-ShareAlike 3.0 %%
%% http://creativecommons.org/licenses/by-sa/3.0/             %%
%%%%%%%%%%%%%%%%%%%%%%%%%%%%%%%%%%%%%%%%%%%%%%%%%%%%%%%%%%%%%%%%


%% Images
\newcommand{\CcImageBy}[1]{%
	\includegraphics[scale=#1]{creative_commons/cc_by_30.pdf}%
}
\newcommand{\CcImageCc}[1]{%
	\includegraphics[scale=#1]{creative_commons/cc_cc_30.pdf}%
}
\newcommand{\CcImageDevNations}[1]{%
	\includegraphics[scale=#1]{creative_commons/cc_dev_nations_30.pdf}%
}
\newcommand{\CcImageNc}[1]{%
	\includegraphics[scale=#1]{creative_commons/cc_nc_30.pdf}%
}
\newcommand{\CcImageNd}[1]{%
	\includegraphics[scale=#1]{creative_commons/cc_nd_30.pdf}%
}
\newcommand{\CcImagePd}[1]{%
	\includegraphics[scale=#1]{creative_commons/cc_pd_30.pdf}%
}
\newcommand{\CcImageSa}[1]{%
	\includegraphics[scale=#1]{creative_commons/cc_sa_30.pdf}%
}
\newcommand{\CcImageSampling}[1]{%
	\includegraphics[scale=#1]{creative_commons/cc_sampling_30.pdf}%
}
\newcommand{\CcImageSamplingPlus}[1]{%
	\includegraphics[scale=#1]{creative_commons/cc_sampling_plus_30.pdf}%
}


%% Groups
\newcommand{\CcGroupBy}[1]{% zoom
	\CcImageBy{#1}%
}
\newcommand{\CcGroupByNc}[2]{% zoom, gap
	\CcImageBy{#1}\hspace*{#2}\CcImageNc{#1}%
}
\newcommand{\CcGroupByNcNd}[2]{% zoom, gap
	\CcImageBy{#1}\hspace*{#2}\CcImageNc{#1}\hspace*{#2}\CcImageNd{#1}%
}
\newcommand{\CcGroupByNcSa}[2]{% zoom, gap
	\CcImageBy{#1}\hspace*{#2}\CcImageNc{#1}\hspace*{#2}\CcImageSa{#1}%
}
\newcommand{\CcGroupByNd}[2]{% zoom, gap
	\CcImageBy{#1}\hspace*{#2}\CcImageNd{#1}%
}
\newcommand{\CcGroupBySa}[2]{% zoom, gap
	\CcImageBy{#1}\hspace*{#2}\CcImageSa{#1}%
}
\newcommand{\CcGroupDevNations}[1]{% zoom
	\CcImageDevNations{#1}%
}
\newcommand{\CcGroupNcSampling}[2]{% zoom, gap
	\CcImageNc{#1}\hspace*{#2}\CcImageSampling{#1}%
}
\newcommand{\CcGroupPd}[1]{% zoom
	\CcImagePd{#1}%
}
\newcommand{\CcGroupSampling}[1]{% zoom
	\CcImageSampling{#1}%
}
\newcommand{\CcGroupSamplingPlus}[1]{% zoom
	\CcImageSamplingPlus{#1}%
}


%********************************************

\begin{document}

% stran 1 med uvodnimi listi
\thispagestyle{empty}

\begin{center}
{\large
UNIVERZA V LJUBLJANI\\
FAKULTETA ZA RAČUNALNIŠTVO IN INFORMATIKO\\
}

\vspace{3cm}
{\LARGE Andraž Vrhovec}\\

\vspace{2cm}
\textsc{\textbf{\LARGE
Varnost podatkov v oblaku
}}

\vspace{2cm}
{ DIPLOMSKO DELO}\\
{ NA INTERDISCIPLINARNEM UNIVERZITETNEM ŠTUDIJU
}


\vspace{2cm}
{\Large Mentor: prof. dr. Aleksandar Jurišić}

\vfill
{\Large Ljubljana, 2016}
\end{center}

\newpage

\ \thispagestyle{empty}

\newpage

%********************************************

% stran 2 med uvodnimi listi
\thispagestyle{empty}

\vspace*{5cm}
{\small \noindent
To diplomsko delo je ponujeno pod licenco \textit{Creative Commons Priznanje avtorstva-Deljenje pod enakimi pogoji 2.5 Slovenija}
ali (po želji) novejšo različico.
To pomeni, da se tako besedilo, slike, grafi in druge sestavine dela kot tudi rezultati diplomskega dela lahko prosto distribuirajo,
reproducirajo, uporabljajo, dajejo v najem, priobčujejo javnosti in predelujejo, pod pogojem, da se jasno in vidno navede avtorja in naslov tega
dela in da se v primeru spremembe, preoblikovanja ali uporabe tega dela v svojem delu, lahko distribuira predelava le pod
licenco, ki je enaka tej.
Podrobnosti licence so dostopne na spletni strani \url{http://creativecommons.si/} ali na Inštitutu za
intelektualno lastnino, Streliška 1, 1000 Ljubljana.

\begin{center}% 0.66 / 0.89 = 0.741573033707865
  \CcImageCc{0.741573033707865}\hspace*{1ex}\CcGroupBySa{1}{1ex}
\end{center}
}

\vspace*{1.5cm}
{\small \noindent
Izvorna koda diplomskega dela, njenih rezultatov in v ta namen razvite programske opreme je ponujena pod GNU General Public License,
različica 3 ali (po želji) novejšo različico. To pomeni, da se lahko prosto uporablja, distribuira in/ali predeluje pod njenimi pogoji.
Podrobnosti licence so dostopne na spletni strani \url{http://www.gnu.org/licenses/}.
}

\begin{center}
\ \\ \vfill
{\em
Besedilo je oblikovano z urejevalnikom besedil \LaTeX.
}
\end{center}

\newpage

\ \thispagestyle{empty}

\newpage

%********************************************

% stran 3 med uvodnimi listi
\thispagestyle{empty}

Namesto te strani {\bf vstavite} original izdane teme diplomskega dela s podpisom mentorja in dekana ter \v zigom fakultete, ki ga diplomant
dvigne v študent\-skem referatu,  preden odda izdelek v vezavo!

\newpage

%********************************************

% stran 4 med uvodnimi listi je prazna
\ \thispagestyle{empty}

\newpage

%********************************************

% stran 5 med uvodnimi listi

\thispagestyle{empty}

\vspace{1cm}
\begin{center}
{\Large \textbf{IZJAVA O AVTORSTVU}}
\end{center}

\begin{center}
{\Large diplomskega dela}
\end{center}

\vspace{1cm}
Spodaj podpisani/-a \hspace{0.5cm} Andraž Vrhovec,

\vspace{0.5cm}
z vpisno številko \hspace{0.5cm} 63080117,

\vspace{1cm}
sem avtor/-ica diplomskega dela z naslovom:

\vspace{0.5cm}
Varnost v oblaku
\vspace{1.5cm}
S svojim podpisom zagotavljam, da:
\begin{itemize}
	\item sem diplomsko delo izdelal/-a samostojno pod mentorstvom

	prof. dr. Aleksandar Jurišić

	\item	so elektronska oblika diplomskega dela, naslov (slov., angl.), povzetek (slov., angl.) ter ključne besede (slov., angl.) identični s tiskano obliko diplomskega dela
	\item soglašam z javno objavo elektronske oblike diplomskega dela v zbirki ''Dela FRI''.
\end{itemize}

\vspace{1cm}
V Ljubljani, dne \today \hspace{1cm} Podpis avtorja/-ice:

\newpage

%*******************************************

% stran 6 med uvodnimi listi je prazna pri dvostranskem tiskanju
\ \thispagestyle{empty}

\newpage

%********************************************

% stran 7 med uvodnimi listi

\chapter*{Zahvala}

\thispagestyle{empty}
ablak
Mogoče :D


\newpage

%********************************************

% stran 8 med uvodnimi listi je prazna pri dvostranskem tiskanju
\ \thispagestyle{empty}

\newpage

%********************************************

% stran 9 med uvodnimi listi
\thispagestyle{empty}

$\;$

\vspace{5cm}
\hfill {\Large \em Morebitno posvetilo}
\thispagestyle{empty}

\newpage

% stran 10 med uvodnimi listi je prazna pri dvostranskem tiskanju

\ \thispagestyle{empty}

\newpage

%********************************************

\renewcommand\thepage{}
\tableofcontents
\renewcommand\thepage{\arabic{page}}

\thispagestyle{empty}


%********************************************

\chapter*{Seznam uporabljenih kratic in simbolov}
\thispagestyle{empty}

\printglossaries

%\cleardoublepage

\clearpage{\pagestyle{empty}\cleardoublepage}

%********************************************
%zacno se glavni listi, ki so numerirani z arabskimi stevilkami

\setcounter{page}{1}
\pagenumbering{arabic}

\chapter*{Povzetek}

\addcontentsline{toc}{chapter}{Povzetek}

\todo[inline]{Napisi povzetek}

\vspace{1.3cm}
\noindent
{\large \bf Ključne besede:}

\vspace{0.5cm}
\noindent
diploma, oblak, varnost, kriptografija


\chapter*{Abstract}

\addcontentsline{toc}{chapter}{Abstract}

Povzetek naj bo napisan v angleškem jeziku.

\vspace{1.3cm}
\noindent
{\large \bf Key words:}

\vspace{0.5cm}
\noindent
Ključne besede v angleškem jeziku.


%********************************************

\chapter{Uvod}

V zadnjih letih smo priča bliskovitemu porastu popularnosti t.i. računalništva v oblaku. Gartner je v začetku leta 2015 uvrstil računalništvo v oblaku med 10 strateško najpomembnejših tehnoloških trendov, ki bodo krojili prihodnost. Oblak nam omogoča udobno, cenovno ugodno delo, kadarkoli in kjerkoli, tudi zunaj pisarne, saj je vse kar potrebujemo internetni dostop in računalnik ali telefon. Računalništvo v oblaku je posledica razvoja več različnih tehnologij, ki so skupaj omogočile cenovno učinkovite rešitve, ki so hkrati preproste za uporabo. Končni uporabniki uporabljajo oblak, ker je preprosto za uporabo, omogoča dostop do podatkov kjerkoli in kadarkoli. Na drugi strani vedno več podjetji v žeji, da znižajo stroške poslovanja in povečajo sposobnost prilagajanja zahtevam trga, seli svoje poslovanje v oblak. Posledica tega premika je selitev podatkov iz fizičnih nosilcev, ki so bili tradicionalno v lasti iste osebe kot podatki, v oblak. To pomeni da varnost hrambe svojih podatkov zaupamo tretjim osebam, kar niti ni nujno tako slabo, kot se sliši. Za veliko ljudi in manjših podjetji je skrb za varnost podatkov velik zalogaj, ki zahteva ustrezno tehnično izobrazbo, strojno opremo in prostor. Kadar hranimo podatke v oblaku, nič od tega ni potrebno, saj nam vse to zagotavlja ponudnik storitev v oblaku v zameno za mesečno naročnino. V nadaljevanju tega dela, bomo najprej pregledali kaj je oblak, kakšne oblike oblaka poznamo. Ko bomo opremljeni z osnovnim poznavanjem oblaka, si bomo pogledali še formalno definicijo varnosti nato pa bomo preučili kako vpeljava in vsesplošna uporaba oblaka vpliva na varnost naših podatkov.

\section{Kaj je oblak?}
Priljubljen internetni rek gre nekako takole: “There is no cloud, just someone else's computer”. Preprosto rečeno je to delno res, saj oblak ni nič drugega kot skupek velikega števila računalnikov, ki niso v lasti uporabnika storitve v oblaku, morda niso niti v lasti ponudnika storitve, ampak jih ta najema pri tretjem ponudniku. Pojem oblak je le krovni izraz za nove načine ponudbe in uporabe storitev, zato ga moramo podrobneje razdelati in razložiti.

NIST opredeli oblak kot: “Računalništvo v oblaku je model, ki omogoča vseprisoten preprost dostop na zahtevo do skupnega računalniških virov (omrežja, strežniki, pomnilniški mediji, aplikacije in storitve), ki so lahko hitro pridobljeni ali sproščeni z minimalnim trudom in popolnoma samodejno, brez posredovanja ponudnika storitve. Model oblaka sestavlja pet ključnih lastnosti, trije nivoji storitve in štirje tipi oblakov.”

\begin{description}[style=nextline]
   \item[Storitev na zahtevo] Stranka lahko pridobi nove vire v oblaku po potrebi, lahko tudi popolnoma samodejno brez človeškega posredovanja.

   \item[Širok dostop preko omrežja] Zmogljivosti oblaka so na voljo preko omrežja in do njih lahko dostopamo preko standardnih mehanizmov, kar omogoča dostops praktično katerekoli naprave ki se je zmožna povezati v omrežje (telefoni, tablice, računalniki, …).

   \item[Uporaba skupnih virov] Ponudnikovi fizični viri so zdrženi v bazen virov iz katerega stranke dobijo virtualizirane vire. Virtualizirani vir lahko prosto prehaja med fizičnimi gostitelji, zato težko govorimo o točni lokaciji vira.

   \item[Bliskovita prožnost] Uporabnik lahko preprosto poveča ali zmanjša zmogljivosti virov, ki jih potrebuje glede na potrebe. Uporabnik ima pogosto občutek da so mu na voljo neomejene zmogljivosti.

   \item[Merjenje porabe in obračun glede na porabo virov]  Oblačni sistemi samodejno nadzirajo in optimizirajo porabo virov s pomočjo metrik (poraba prostora, pasovne širine, procesorskega časa …). Poraba virov se nadzira, kontrolira in poroča, kar omogoča transparentnost za ponudnika in uporabnika.
\end{description}

\newacronym{SaaS}{SaaS}{Software as a Service}
\newacronym{PaaS}{PaaS}{Platform as a Service}
\newacronym{IaaS}{IaaS}{Infrastructure as a Service}
\newacronym{API}{API}{Application programming interface}

\begin{description}[style=nextline]
	\item[\gls{SaaS}] Ponudnik storitev uporabniku omogoča uporabo njegove aplikacije, ki teče v oblaku. Aplikacija je dosegljiva iz različnih naprav, preko recimo spletnega brskalnika ali aplikacijskega vmesnika (\gls{API}). Uporabnik nima nadzora nad infrastrukturo na kateri teče aplikacija (omrežje, strežniki …), kot tudi ne nad samo aplikacijo (z izjemo konfiguracije aplikacije).

	\item[\gls{PaaS}]: Ponudnik storitev uporabniku ponuja platformo, na katero lahko uporabnik namesti lastne ali kupljene aplikacije, ki so prilagojene za delovanje v oblaku. Tudi tu uporabnik nima nadzora nad infrastrukturo na kateri teče aplikacije, ima pa popolni nadzor nad aplikacijo.

	\item[\gls{IaaS}]: Ponudnik uporabniku omogoča postavitev osnovnih gradnikov računalniškega sistema (procesiranje, hramba podatkov, omrežje …) iz katerih lahko uporabnik sestavi kar želi. Uporabnik sicer nima nadzora nad spodaj ležečo fizično infrastrukturo, vendar ima ves nadzor nad virtualiziranimi viri ki jih pridobi.
\end{description}

\begin{description}[style=nextline]
	\item[Zasebni oblak] Vsa infrastruktura je namenjena uporabi znotraj ene organizacije, ki je sestavljena iz več delov. Upravljanje, vzdrževanje in lastništvo je lahko v domeni iste organizacije, tretje osebe ali kombinacije obojega. Lokacija sistema je lahko v prostorih organizacije ali pa izven.

	\item[Skupni oblak] Infrastruktura je namenjena uporabi s strani skupnosti organizacij, ki imajo skupne zahteve ali skupni cilj. Upravljanje, vzdrževanje in lastništvo je lahko v domeni ene ali večih organizacij znotraj skupnosti, tretje osebe ali kombinacije obojega. Lokacija sistema je lahko v prostorih ene od organizacij ali pa izven.

	\item[Javni oblak] Infrastruktura je namenjena uporabi s strani vseh. Ponudnik je lahko zasebno podjetje, izobraževalna ustanova ali vladna agencija. Infrastruktura se ponavadi nahaja na lokaciji ponudnika.
	Hibridni oblak: Hibridni oblak je kombinacija dveh ali več tipov oblakov iz prvih treh točk (zasebni, javni, skupni). Oblaka sta sicer fizično ločena, vendar skupni standardi in tehnologije omogočajo preprosto prehajanje aplikacij in podatkov iz enega na drugega.
\end{description}

\section{Kaj je varnost?}
Varnost je širok pojem, ki ga v kontekstu računalništva lahko razdelimo na tri dele. To so: računalniška varnost (computer security), informacijska varnost (information security) in varovanje podatkov (information assurance).

\newglossaryentry{neizpodbitnost}
{
  name=preprečevanje tajenja,
  description={non-repudiation},
}

\begin{description}

	\item[Računalniška varnost] se ukvarja z varovanjem računalniških sistemov, od strojne opreme preko programske opreme do podatkov, ki se procesirajo na tej strojni opremi. Naloga računalniške varnosti je zagotoviti integriteto, dostopnost in zaupnost vseh sestavnih delov računalniškega sistema.

\item[Informacijska varnost] se ukvarja z varovanjem podatkov, ki se obdelujejo in shranjujejo v informacijskih sistemih. Njena naloga je zasčititi podatke pred nepooblaščenim dostopom, uporabo, spreminjanjem in uničenjem z ciljem zagotavlanja intgritete, dostopnosti in zaupnosti podatkov.

\item[Varovanje podatkov] se ukvarja z varovanjem podatkov in upravljanjem tveganj povezanih z uporabo, prenosom in procesiranjem podatkov z ciljem zagotavljanja integritete, dostopnosti, zaupnosti in neizpodbitnostjo podatkov.
\end{description}

Iz teh opisov lahko razberemo kljucne lastnosti varnega računalniškega sistema. V literaturi se te ključne tri lastnosti imenujemo CIA triad, kjer je CIA angleška kratica za ključne 3 lastnosti:

\newglossaryentry{zaupnost}
{
  name=zaupnost,
  description={Confidentiality},
}

\newglossaryentry{celovitost}
{
  name=celovitost,
  description={Integrity},
}

\newglossaryentry{razpolozljivost}
{
  name=razpoložljivost,
  description={Availability},
}

\begin{description}
\item[\Gls{zaupnost}] pomeni zasebnost podatkov. Ukrepi, ki zagotavljajo zaupnost podatkov so usmerjeni k preprečevanju dostopa nepooblaščenih oseb. Pogosti mehanizmi, ki se uporabljajo za zagotavljanje zaupnosti podatkov so zaščita z uporabniškim imenom in geslom, dvostopenjska overitev (two-factor auth), biometrična overitev in enkripcija podatkov.

\item[\Gls{celovitost}] pomeni ohranjanje točnosti in pravilnosti podatkov. Ukrepi, ki zagotavljajo celovitost podatkov, preprečujejo nepooblaščenim osebam spreminjanje podatkov med prenosom ali v mirovanju. Omogočajo sledenje spremembam podatkov skozi zgodovino (revzijska sled), obnovitev podatkov iz varnostne kopije v primeru da pride do nepooblaščene spremembe. Mehanizmi, ki zagotavljajo celovitost podatkov so nadzor dostopa in zgoščevalne funkcije in digitalni podpisi.

\item[\Gls{razpolozljivost}] pomeni da so podatki vedno na voljo pooblaščenim osebam. Glavni vzrok nerazpoložljivosti so programske in strojne napake ter napadi zavrnitve storitve (denial of service). Razpoložljivost podatkov se zagotavlja z rednim vzdrževanjem strojne in programske opreme in s posebnimi ukrepi za preprečevanje napadov zavrnitve storitve.

\end{description}

Poleg treh ključnih lastnosti, ki jih najdemo zgoraj je Meier identificiral še tri lastnosti, to so:

\newglossaryentry{overjanje}
{
  name=overjanje,
  description={Authentication},
}

\newglossaryentry{avtorizacija}
{
  name=avtorizacija,
  description={Authorization},
}

\newglossaryentry{neizpodbitnost}
{
  name=neizpodbitnost,
  description={Non-repudioation},
}

\begin{description}

\item[\Gls{overjanje}] identificira vir podatka oz. izvor operacije, ki je lahko človek ali računalnik in se prepričamo da je res tisti za katerega se izdaja.

\item[\Gls{avtorizacija}] nadzira dostop do virov glede na pravila sistema in dodeljene pravice. Avtorizacijski mehanizmi uporabnikom preprečujejo da bi presegli svoja pooblastila.

\item[\Gls{neizpodbitnost}] onemogoča uporabniku da izpodbija oz. zanika neko dejanje ali veljavnost pogodbe.

\end{description}

Vse zgoraj naštete lastnosti so splošne lastnosti varnosti in se kot take seveda tičejo tudi računalništva v oblaku. Pri uporabi storitev v oblaku bi morali najbolj izpostaviti vprašanje varnosti podatkov, saj le te zaupamo v hrambo ponudniku storitve.

\chapter{Varnostna tveganja v oblaku}

Zaradi vse večje popularnosti storitev v oblaku in posledično ogromne količine podatkov, s katerimi razpolagajo ponudniki storitev v oblaku, so ti vedno pogosteje tarča napadov kriminalnih združb. Za sisteme, ki tečejo “v oblaku”, je značilno da navzven delujejo kot en strežnik, v resnici pa tečejo na velikem številu strežnikov, ki so med seboj povezani. To pomeni da morebitni vdor v samo en strežnik napadalcu omogoča dostop tudi na druge strežnike, ki tečejo znotraj istega sistema. Če si sistem predstavljamo kot trdnjavo, to pomeni da je večina zaščitnih mehanizmov skoncentrirana na obrobju, na obzidju trdnjave. Če uspemo to obzidje predreti na samo enem delčku, se potem lahko več ali manj prosto sprehajamo znotraj trdnjave. Pred pojavom storitev v oblaku je bilo znotraj enega obzidja le nekaj strežnikov, zato je bilo obzidje malo in ga je bilo lažje vzdrževati. Sistem, na katerem teče storitev v oblaku pa je lahko sestavljen iz več tisoč strežnikov, kar pomeni da velikost obzidja močno poveča in veliko obzidje je težko vzdrževati, zato se kaj hitro v njem pojavijo razpoke.

Druga stvar, ki dela sisteme v oblaku zanimive za napadalce je koncentracija podatkov. V primeru do vdora v klasični strežnik, napadalec dobi dostop do veliko manjše količine podatkov kot ob vdoru v oblačni sistem. Razlog je preprost, saj na en sam strežnik fizično ni mogoče shraniti toliko podatkov, kot na sistem, ki je sestavljen iz tičoče strežnikov.

Ranljivost računalniškega sistema je lastnost, zaradi katere je ogroženo njegovo delovanje in varnost podatkov, ki jih vsebuje. Storitve, ki delujejo v oblaku tako podedujejo vse ranljivosti, ki jih imajo klasični računalniški sistemi, zaradi posebnih lastnosti oblaka, pa obstaja se nekaj groženj, specifičnih za storitve v oblaku:

\begin{description}[style=nextline]

\item[Nedovoljen dostop do upraviteljskega in programskega vmesnika] Narava storitev v oblaku zahteva da so upraviteljski vmesniki dostopni preko Interneta siroki množici uporabnikov, kar posledično pomeni da so ti vmesniki izpostavljeni napadom preko Interneta.

\item[Ranljivosti internetnega protokola] Ena ključnih lastnosti oblaka je tudi široka dostopnost preko standardiziranih protokolov kot je Internetni protokol, zato je potrebno ko govorimo o varnosti v oblaku upoštevati tudi vse ranljivosti Internetnega protokola, kot so recimo napadi vrinjenega napadalca (Man in the middle attack).

\item[Ranljivosti obnovitve podatkov] Zaradi uporabe skupnih virov, ki je ena ključnih lastnosti oblaka, se lahko en fizični nosilec podatkov, ki je prej pripadal enemu uporabniku kasneje dodeli drugemu uporabniku. Če pred tem ne poskrbimo za varen izbris podatkov, lahko drugi uporabnik obnovi izbrisane podatke z nosilca in tako pride do razkritja podatkov.

\item[Izogibanje merjenju in plačilu] Storitve v oblaku se običajno plačujejo po porabi. Prav tako, da pričnemo uporabljati storitev praviloma potrebujemo vnesti samo nekaj osnovnih podatkov in številko kreditne kartice. Napadalec nam lahko predloži ukradene podatke ali manipulira podatke o porabi, tako da se izogne plačilu.

\end{description}

\subsection{Izzivi na področju varnosti v oblaku}

\subsubsection{Razkritje podatkov}
O razkritju podatkov govorimo v primeru, da se občutljivi in zaupni podatki znajdejo na očeh nepooblaščenih oseb ali javnosti.  To se lahko zgodi kot posledica ciljanega napada ali preprosto zaradi človeške napake, napake v aplikaciji ali slabih varnostnih mehanizmov aplikacije. Primeri zaupnih podatkov, ki so pogosto tarča napadov so: zdravstveni podatki, finančni podatki, osebno določljivi podatki, poslovne skrivnosti in intelektualna lastnina.

Razkritje podatkov je stara in vseprisotna grožnja v vseh računalniških sistemih, vendar zaradi prej omenjene koncentracije podatkov je ta grožnja v oblačnih sistemih še večja.

Obseg škode, ki ga povzroči razkritje podatkov je odvisen od občutljivosti podatkov, ki so bili odtujeni. V veliko državah po svetu obstajajo zakoni in predpisi, ki podjetja ki hranijo določene tipe občutljivih podatkov obvezuje k določenim standardom zaščite. V primeru, da pride do razkritja podatkov, lahko odgovorno podjetje doleti denarna kazen ali celo kazenska ovadba. Poleg škode, ki je bila povzročena z razkritjem podatkov in morebitno denarno kaznijo so tu še stroški preiskave incidenta, obveščanja uporabnikov in izguba ugleda podjetja.

\subsubsection{Pomankljivo upravljanje identitete, poverilnic in nadzora dostopa}
Površno upravljanje lahko pripelje do nenamernega razkritja uporabniških imen, gesel in ključev, ki se uporabljajo za dostop do storitev. Ponudniki storitev v oblaku upravljajo z identitetami milijonov uporabnikov in zaposlenimi, zato je potrebna velika skrb da se zagotovin pravilno upravlanje identitete skozi celoten cikel. Potrebno je poskrbeti za takojšno prekinitev veljavnosti identitete, v primeru da pride do odpovedi delovenega razmerja zaposlenega ali prekinitve pogodbe s stranko. Uporaba šibkih gesel in neuporaba dvostopenjske avtentikacije in neuporaba rotacije šifrirnih ključev, gesel in certifikatov so primeri slabe prakse, ki v primeru napada napadalcu olašajo delo in povečajo razsežnosti škode, ki jo napadalec lahko povzroči.

\subsubsection{Nezaščiteni uporabniški in programski vmesniki}
Ponudniki storitev v oblaku ponujajo nabor uporabniški in programskih vmesnikov, preko katerih uporabniki dostopajo do storitev ponudnika. So stična točka med uporabnikom in storitvijo in od varnosti in dostopnosti teh vmesnikov je odvisna varnost celotnega sistema. So prva linija obrambe sistema, zato je pri zasnovi teh vmesnikov potrebno misliti na varnost s pravilno implementacijo varnostnih mehanizmov kot so avtentikacija, nadzor dostopa in enkripcija. Vmesniki morajo biti zasnovanai tako, da ščitijo pred namerno in nenamerno zlorabo sistema.

\subsubsection{Sistemske ranljivosti}
Sistemske ranljivosti so preprosto povedano hrošči v uporabljeni programski opremi, katere je možno izkoristiti za napad na računalniški sistem. Ranljivosti, ki se nahajajo v jedru opreacijskega sistema ali sistemskih knjižnicah predstavljajo grožnjo za celoten sistem.

Tako kot večina ostalih predstavljenih groženj, tudi sistemske ranljivosti niso nič novega, vendar razsežnost oblačnih sistemov doda takim ranljivostim popolnoma nov pomen. Zaradi napredka tehnologije je možno na enem strežniku gostiti več virtualnih strežnikov različnih organizacij, kar pomeni da si ti virtualni strežniki delijo pomnilnik in ostale vire, kar omogoča nove načine izrabe sistemskih ranljivosti.

\subsubsection{Ugrabitev uporabniškega računa}
Ugrabitve uporabniških računov so posledica napadov ribarjenja, prevare ali izkoriščanja sistemskih pomankljivosti. V primeru, da ista uporabniška imena in gesla uporabljamo za več različnih storitev, to še dodatno poveča posledice napada. V primeru, da pride do ugrabitve uporabnišlega računa za storitev v oblaku, lahko napadalec izkoristi to storitev kot odskočno desko za nove napade.

\subsubsection{Notranje grožnje}
Tveganje, ki ga predstavljajo notranje grožnje, kot so škodoželjni zaposleni, je dokaj novo in še vedno predmet debate v industriji računalniške varnosti. Dejstvo je da tako tveganje obstaja, predvsem pri uporabi javnega oblaka, saj v tem primeru z infrastrukturo in podatki razpolaga tretja oseba. V tem primeru uporabnik storitve nima nikakršnega mehanizma, s katerim bi lahko nadziral zaposlene pri ponudniku storitve.

\subsubsection{Organizirana trajna grožnja}
Organiziranje trajne grožnje so ena najnaprednejših oblik grožnje računalniški varnosti. Za napadi te oblike so ponadavi skrivajo velike organizacije z državno podporo (tajne službe, vojska). Značilnost te oblike napada je, da napadalec temeljito preuči ciljni sistem, njegovo organizacijo in šibke točke. Ta oblika napada se pogosto začne z usmerjenim zvablanjem, to je posebna tehnika ribarjenja, ki je sumerjena proti točno določenim osebam, v našem primeru zaposleni v ciljni organizaciji. S pomočjo usmerjenega zavblanja napadalec poskuša pridobiti vstopno točko v ciljni sistem. Ko enkrat pridobi dostop do notranjega omrežja sistema, to uporabi kot odskočno desko za nadalnje premike proti željenim tarčam. Cilj napada so podatki, ki se nahajajo v sistemu žrtve, pogosto so to poslovne skrivnosti ali državne tajne. Napadalec želi čimdlje ostati skrit, saj tako pridobi dostop do večjega števila računalnikov znotraj omrežja in s tem večje količine podatkov. V želji da ostane skrit, uporablja napredne tehnike brisanja sledov, zato je take tipe napada zelo težko zaznati, saj navzven sistem izgleda kot da deluje normalno.

Posebna oblika organiziranih trajnih groženj so kibernetska orožja. Verjetno najbolj znan med njimi je Stuxnet, ki js sabotiral Iranski jedrski program. Kibernetska orožja se od ostalih oblik napadov razlikujejo v tem, da žrtvi povzročijo dejansko fizično škodo v obliki uničene opreme.

\subsubsection{Izguba podatkov}
Izguba podatkov predstavlja za večino uporabnikov in podjetji eno najhujših nočnih mor. Trajna izguba podatkov, shranjenih v oblaku ni nujno posledica napada tretje osebe, ampak je lahko tudi posledica malomarnosti, kot je recimo nenameren izbris podatkov s strani ponudnika storitve ali posledica naravne katastrofe kot je požar ali potres. V izogib trajni izgubi podatkov ponudniki storitev v oblaku uporabljajo različne mehanizme, kot so gegografska redundanca in varnostno kopiranje na drugo lokacijo. Vendar v primeru, da uporabnik podatke kriptira še preden jih posšlje v oblak in nato izgibu enkripcijski ključ, noben od naštetih mehanizmov ne more preprečiti takojšne in trajne izgube podatkov.

\subsubsection{Malomarnost}
Preden začnemo uporabljati storitve v oblaku, je potrebno preučiti vse aspekte, ki jih bo taka sprememba prinesla. Vodilni kadri v podjetjih pogosto nimajo dovolj poglobljenega poznavanja oblaka s tehnološkega vidika, kot tudi ne poznajo pravnih zank, ki jih vsebujejo “pogoji uporabe” določenih ponudnikov storitev v oblaku. V primeru, da pred uporabo storitev v oblaku ne naredimo temeljite analize, lahko pride do tehnoloških, pravnih, skladenjskih in poslovnih tveganj, ki ogrozijo uspeh podjetja.

\subsubsection{Zloraba in zlonamerna uporaba sistemov}
Pogosto ponudniki storitev v oblaku omogočajo brezplačne preizkusne račune za omejeno časovno obdobje, pa tudi če jih ne, lahko napadalec s pomočjo ukradenih bančnih podatkov ustvari veliko količino uporabniških računov, ki jih potem izkoristi za izvajanje različnih nelegalnih aktivnosti. Te lažni uporabniški računi napadalcu omogočajo pošiljanje nezaželjene pošte, izvajanje napadov ribarjenja ali napadov zavrnitve storitve, ki so lahko usmerjeni k ponudniku storitve, njegovim strankam ali pa tretjim osebam.

\subsubsection{Napadi zavrnitve storitve}
Napad zavrnitve storitve je ena najpreprostejših in pogosto zelo učinkovitih oblik napada na računalniške sisteme. Cilj napada je uporabnikom preprečiti normalen dostop do storitve, kar napadalec doseže s tem da napaden sistem prisili v nerazumno veliko porabo različnih virov, kot so pasovna širina, procesorski čas ali podatkovni prostor. Najbol osnovna oblika napada zavrnitve storitve se zanaša na dejstvo da napadalec razpolaga z večjo količino določenega vira (pasovna širina, procesorski čas …) kot napaden sistem. Obstajajo pa tudi bolj prefinjene asimetrične oblike napada, kjer napadalec z veliko manjšo količino virov kot tarča, s pomočjo izrabe varnostne pomankljivosti v sistemu ali aplikaciji onemogoči tarčo napada. Napadi zavrnitve storitve običajno ne pustijo trajnih posledic, vendar predstavljajo veliko neprijetnost za uporabnike, saj le ti za čas trajanja napada ne morejo uporabljati storitve, prav tako pa škodujejo ugledu ponudnika.

\subsubsection{Ranljivosti zaradi uporabe skupnih virov}
Zadnji sklop varnostnih groženj v oblaku predstavljajo ranljivosti, ki so posledica ene od ključnih lastnosti oblaka, to je uporaba skupnih virov. Ponudniki storitev v oblaku zagotavljajo navidezno neomejene vire končnemu uporabniku s pomočjo tehnologij virtualizacije in skupne uporabe. Pri tem uporabljajo klasično strojno opremo, ki sicer s pomočjo posebne programske opreme omogoča delitev virov, vendar ni bila zasnovana, da bi omogočala varno skupno uporabo. Tak tip napadov izkorišča dejstvo, da si različni virtualni strežniki delijo fizično strojno opremo, zato je možno s posebno obliko napada, ki se imenuje napad preko stranskega kanala pridobiti informacije iz žrtvinega strežnika, brez da bi dejansko prišlo do vdora. V času pisanja še ni prišlo do takega tipa vdora v praksi, vendar obstajajo primeri takih napadov, ki so jih demonstrirali pod laboratorijskimi pogoji.










\newglossaryentry{cistopis}
{
  name=čistopis,
  description={plaintext},
}


\newglossaryentry{tajnopis}
{
  name=tajnopis,
  description={ciphertext},
}

\newglossaryentry{sifrirnikljuc}
{
  name=šifrirni ključ,
  description={encryption key},
}


\chapter{Zaščita podatkov s pomočjo kriptografije}

Kriptografija omgoča doseganje informacijske varnosti s pomočjo matematičnih postopkov. Verjetno najbolj poznan tak postopek je šifriranje, ki zagotavlja zaupnost podatkov, pogled tega v področje kriptografije spadajo tudi digitalni podpisi, avtentikacijske kode MAC in zgoščevalne funkcije, ki nam pomagajo pri ohranjanju celovitosti in overjanju identitete podatkov.

Šifriranje je način, kako skriti vsebino sporočila (\gls{cistopis}a) pred nepooblaščenimi očmi. To se doseže z uporabo algoritma, ki mu kot vhodni podatek damo \gls{cistopis} in \gls{sifrirnikljuc}, on pa nam vrne \gls{tajnopis}. \Gls{tajnopis} vsebuje enake informacije kot \gls{cistopis}, vendar ima to lastnost, da ga je mogoče prebrati le v primeru, da posedujete ključ za dešifriranje. V kolikor takega ključa ne posedujete, je iz \gls{tajnopis}a praktično nemogoče razbrati vsebino originalnega sporočila.\footnote{Težanost je odvisna od moči kriptografske sheme in parvilnosti implementacije}

To pomeni, da lahko tajnopis javno objavimo, varovati je potrebno samo ključ, s katerim se le ta odšifrira. Ta lastnost nam omogoča, da varno poslujemo s svojo banko, pošiljamo e-pošto, ki jo lahko prebere samo prejemnik navkljub temu, da Internet ni varen komunikacijski kanal in na njemu mrgoli zlobnežev, ki želijo priti do naših zaupnih podatkov.

\

Da nam bo lažje, najprej definirajmo kaj je kriptografska shema ali kriptosistem.

\begin{mdframed}[frametitle={Kriptografska shema}]
Kriptografska shema je peterka  $(\mathcal{P}, \mathcal{C}, \mathcal{K},  \mathcal{E}, \mathcal{D})$, kjer veljajo naslednje lastnosti:
	\begin{enumerate}
		\item $\mathcal{P}$ je končna množica možnih čistopisov
		\item $\mathcal{C}$ je končna množica možnih tajnopisov
		\item $\mathcal{K}$ je končna množica možnih ključev
		\item Ta vsak $K \in \mathcal{K}$ obstaja šifrirni postopek $e_K \in \mathcal{E}$ in pripadajoč dešifrirni postopek $d_K \in \mathcal{D}$. Za funkciji $e_K: \mathcal{P} \rightarrow \mathcal{C}$ in $d_K: \mathcal{C} \rightarrow \mathcal{P}$ velja $d_k(e_k(x)) = x$ za vsak $x \in \mathcal{P}$.
	\end{enumerate}
\end{mdframed}

Zgornjo matematično definicijo se da povzeti z naslednjimi besedami. Kriptografsko shemo sestavlja pet sestavin. Čistopisi, tajnopisi, ključi, šifrirni in dešifrirni postopek. Pomembno je, da ob danem ključu in šifrirnem postopu, obstaja dešifrirni postopek, ki nam vsak tajnopis, ki je bil šifriran z izbranim šifrirnim postopkom nedvoumno pretvori nazaj v čistopis.

Skozi zgodovino človeštva se je že zelo zgodaj pojavila želja po skrivanju sporočil in s tem so se razvile tudi prve metode šifriranja, vendar nobena od njih ni prestala preiskusa časa.

\

Šifrirne algoritme delimo na dve skupini:
\begin{enumerate}

	\item Simetrične oz. klasične šifre
	\item Asimetrične šifre oz. kriptografija javnih ključev

\end{enumerate}

Simetrične šifre so najstarejša in najbolj razširjena vrsta šifriranja, kamor spadajo tudi vse zgodovinske šifre, ki se ne uporabljajo več. Poglavitna lastnost simetričnih šifer je, da se za šifriranje in dešifriranje uporablja enak ključ. Take šifre so lahko zelo preproste in hitre, vendar se pojavi problem distribucije ključev. Tajnopise lahko prenašamo po nezavarovanih kanalih, kar pa ne velja za ključe.

Ta problem lahko rešijo asimetrične šifre, ki uporabljajo en ključ za šifriranje, drug ključ pa dešifriranje. Tak par ključev imenujemo javni in zasebni ključ. Kot že ime pove, javni ključ ni nobena skrivnost in ga lahko brez skrbi objavimo na javnem mestu, obratno pa velja za zasebni ključ. Sporočilo, šifrirano z javnim ključem je možno dešifrirati samo s pripadajočim zasebnim ključem.

\subsection{Simetrične šifre}

V kategorijo simetričnih šifer spadajo vse šifre, pri katerih se za šifriranje in dešifriranje uporablja enak ključ ali pa obstaja bijektivna preslikava, ki nam šifrirni ključ preslika v dešifrirni in obratno. Simetrične šifre še naprej delimo na tokovne in bločne šifre. Poglavitna razlika med tokovnimi in bločnimi šiframi je v načinu šifriranja podatkov.

Naj bo dano sporočilo $X$ velikosti $n$ ($|X|=n$), ki je sestavljeno iz $n$ osnovnih nosilcev informacij (to so lahko biti, bajti ali kaj čisto tretjega). Simobol $\|$ predstavlja združevanje (lepljenje) dveh nizov.
$$
X=x_0\|x_1\|x_2\|...\|x_{n-1}
$$

Bločna šifra vzame blok (del sporočila) velikosti $m$; $n=im, i\in\mathbb{N} $ in ga šifrira s ključem $K$. Tajnopis $Y$ tako pridobimo po naslednjem postopku:
\begin{align*}
1*1
\end{align*}


\begin{gather*}
	y_j = e_K(e_{j} \| e_{j+1} \| \ldots \| e_{j+m-1})  \quad j=0,1,\ldots,i\\
	Y = y_0\|y_1\|\ldots\|y_i
\end{gather*}

Bločne šifre delujejo na bloku podatkov fiksne velikosti (običajno velikost ključa določi velikost bloka) in v kolikor dolžina sporočila ni deljiva z $m$, je potrebno sporočilo doploniti s posebnim zaporedjem znakov, dokler dolžina ni večkratnik števila $m$.

Drugačen pristop pa uporabljajo tokovne šifre, ki iz ključa $K$ generirajo neskončno zaporedje ključev $z$, ki ga nato uporabijo za šifriranje sporočila.

\begin{gather*}
z=z_0\|z_1\|z_2\|\ldots\|z_n \\
y_j=e_{z_j}(x_j) \quad j=0,1,\ldots,n\\
Y = y_0\|y_1\|\ldots\|y_n \\
\end{gather*}

Tak način šifriranja lahko uporabimo za poljubno dolga sporočila, brez potrebe po dopolnjevanju.

\newacronym{DES}{DES}{Data encryption standard}
\newacronym{NIST}{NIST}{National Institute of Standards and Technology}
\newacronym{NBS}{NBS}{National Bureau of Standards}


\newglossaryentry{feistelovasifra}
{
  name=Feistlova šifra,
  description={encryption key},
}


\subsubsection{DES in 3DES}

\gls{DES} je bil eden prvih modernih šifrirnih algoritmov. Originalna implementacija je nosila ime Lucifer in bila razvita v sedemdesetih letih v laboratorijih podjetja IBM. Kasneje je bila rahlo modificirana različica tega algoritma sprejeta s strani \acrshort{NIST}-a (takrat se je imenoval še \acrshort{NBS}) kot standard \gls{DES}.

\gls{DES} je simetrična bločna šifra, ki deluje na blokih velikosti 64 bitov in uporablja 64 bitni ključ, vendar je v ključu možno poljubno izbrati le 56 bitov, preostalih 8 pa se uporablja samo za odkrivanje napak v ključu.

Algoritem je sestavljen iz 16 identičnih korakov oz. krogov, nekaj pre- in post-procesiranja vhodnih podatkov, Feistelove F-funkcije ter funkcije, ki iz glavnega ključa generira zaporedje podključev. Na začetku se 64-bitni blok razdeli na dve 32-bitni polovici, levo in desno, ki se po vsakem končanem krogu algoritma zamenjata. Taka shema delovanja je znana tudi kot \gls{feistelovasifra}. V vsakem krogu se leva polovica bloka z opreacijo $XOR$ zmeša z rezultatom Feistelove F-funkcije, ki kot vhodne parametre dobi podključ trenutnega kroga in desno polovico bloka. Po koncu vsakega kroga se leva in desna polovica zamenjata.

Ob prelomu tisočletja se je izkazalo da je \gls{DES}-ov 64-bitni o. 56-bitni ključ prešibak, da bi zadržal napade s surovo silo. Zato se je pojavila varianta algoritma imenovana 3DES. 3DES uporablja 3 ključe, pri čemer imamo na voljo 3 opcije za izbiro ključev.

\begin{enumerate}
\item Vse tri ključe izberemo neodvisno
\item Izberemo ključa $K_1$ in $K_2$ in $K_3=K_1$
\item Vsi ključi so isti, torej $K_1=K_2=K_3$
\end{enumerate}

Postopek šifriranja in dešifriranja je nato tak:

\begin{gather*}
C=E_{K_3}(D_{K_2}(E_{K_1}(P))) \\
P=D_{K_1}(E_{K_2}(D_{K_3}(C)))
\end{gather*}

Kjer je $C$ tajnopis in $P$ čistopis.

\newacronym{AES}{AES}{Advanced encryption standard}


\subsubsection{AES}

\gls{AES} je bil leta 2001 izbran kot naslednik standarda \gls{DES}. Med 15 prijavljenimi algoritmi je \gls{NIST} izbral družino algoritmov Rijndael. Kot standard \gls{AES} so bile sprejete tri različice algoritma. Vse tri uporabljajo 128-bitne bloke, razlikujejo pa se v velikosti ključa. Na voljo so 3 opcije: 128, 192 ali 256-bitni ključ.

\gls{AES} se po strukturi algoritma razlikuje od \gls{DES}-a, saj namesto Feistelove šifre uporablja substitucijsko-permutacijsko omrežje.  Algoritem iz pričje persepktive izgleda tako:

\begin{description}
	\item[KeyExpansions] S posebnim algoritmom se na podlagi glavnega ključa izračunajo 128-bitni podključi, ki se bodo uporabili v vsakem krogu šifriranja.
	\item[AddRoundKey]  S pomočjo operacije $XOR$ se bloku čistopisa primeša 128-bitni podključ.
	\item[SubBytes] Substitucija bajtov v bloku po pravilu ki ga določa S-škatla.
	\item[ShiftRows] Permutacija bajtov v bloku
	\item[MixColumns] Linearna transformacija
\end{description}

\begin{description}[style=nextline]
	\item[KeyExpansions]

	\item[Prvi krog]
		\begin{enumerate}
			\item AddRoundKey
		\end{enumerate}
	\item[Vmesni krogi]
		\begin{enumerate}
			\item SubBytes
			\item ShiftRows
			\item MixColumns
			\item AddRoundKey
		\end{enumerate}

	\item[Zadnji krogi]
		\begin{enumerate}
			\item SubBytes
			\item ShiftRows
			\item AddRoundKey
		\end{enumerate}

\end{description}

\subsection{Asimetrične šifre}

Eden večjih problemov simetričnih šifrirnih algoritmov je distribucija ključev. V kolikor želita dve osebi svojo komunikacijo zaščititi z uporabo simetričnega šifrirnega algoritma, se morata najprej izmenjati ključ. Dobra praksa je tudi, da se izognemo ponovni uporabi istega ključa, kar v praksi pomeni, da si moramo z vsako osebo, s katero želimo komunicirati izmenjati ključ in hitro se izkaže, da omrežje $n$ ljudi potrebuje $\frac{n(n-1)}{2}$ ključev. Če to preslikamo na omrežje kot je Internet, ki ga uporablja približno 40 odstotkov svetovne populacije. Če vzamemo pesimistično oceno da si vsak uporabnik Interneta lasti samo eno elektronsko napravo, ki uporablja varno komunikacijo je, je naš $n \approx 3.500.000.000$.

Rešitev za ta problem so lahko asimetrični šifrirni algoritmi, oziroma kriptografija javnih ključev. Asimetrični algoritmi ne uporabljajo samo enega, ampak dva ključa. Javni ključ poznajo vsi in ga uporabijo za šifriranje sporočil. Sporočila, šifrirana z javnim ključem, lahko dešifrira samo imetnik zasebnega ključa.

Asimetrični algoritmi delujejo na predpostavki, da so nekatere matematični v eno smer preprosto izračunljivi, v drugo pa preprosto pretežki, da bi jih bilo možno rešiti v razumnem časovnem okvirju. Eden takih problemov je faktorizacija števil. Če si izmislite dve dokaj veliki praštevili $p$ in $q$, ju med seboj pomnožite, dobite novo število $n$, za katero ni tako preprosto ugotoviti iz katerih faktorjev je sestavljeno, ob predpostavki, da ne poznamo $p$ in $q$.

Na žalost pa so zaradi svoje povezanosti s težkimi matematičnimi problemi pogosto asimetrični šifrirni algoritmi veliko počasnejši od svojih simetričnih sorodnikov. V praksi se uporablja kombinacija asimetričnih in simetričnih algoritmov. Asimetrični algoritmi so sicer počasnejši, vendar veliko bolj praktični zaradi preproste izmenjave ključev.


V začetni fazi vzpostavljanja varnega komunikacijska kanala se zaradi preprostejše distribucije ključev uporabi asimetrični algoritem. Ko je varni kanal vzpostavljen se preko njega izmenja sejni ključ, to je ključ ki se

\subsubsection{Izmenjava ključa Diffe-Hellman}

Metoda za izmenjavo ključa Diffe-Hellman je bil prvi javno objavljen algoritem, ki uporablja princip javnih in zasebnih ključev. V praksi omogoča, da se dva neznanca preko nezavarovanega kanala kot je Internet dogovorita za skupni ključ, ki ga bosta uporabljala za nadaljnjo komunikacijo. Prva sta ga objavila Whitfield Diffe in Martin Hellman leta 1976, vendar se je kasneje izkazalo, da je angleška tajna služba odkrila princip kriptografije javnih ključev že leto poprej, vendar to ostalo tajno do leta 1997.

Metoda temelji na reševanju problema diskretnega logaritma, za katerega matematiki verjamejo, da ne obstaja algoritem, ki ga bi bil sposoben rešiti v sprejemljivih časovnih okvirjih (izjema so algoritmi, ki tečejo na kvantnih računalnikih, vendar je to povsem drug svet).

Originalna različica algoritma uporablja multplikativno grupo števil po modulu $p$, kjer je $p$ praštevilo in $g$ primitivni koren števila $p$.

Recimo da se želita Anita in Bojan preko Interneta dogovoriti o skupnem ključu, ki ga bosta uporabljala za nadaljnjo komunikacijo. To lahko storita po sledečem postopku:

\begin{mdframed}[frametitle={Izmenjava ključa Diffe-Hellman}]
\begin{enumerate}
	\item Anita in Bojan se dogovorita os skupnih parametrih, ki jih bosta uporabljala v algoritmu. To je praštevilo $p$ in primitivni koren praštevila $g$
	\item Anita si izbere skrivno število $a$, izračuna $A=g^a \mod p$ in pošlje $A$ Bojanu.
	\item Bojan si izbere skrivno število $b$, izračuna $B=g^b \mod p$ in pošlje $B$ Aniti.
	\item Anita izračuna $s=B^a \mod p$
	\item Bojan izračuna $s=A^b \mod p$
	\item Anita in Bojan sta dobila enak rezultat $s$, ki ga sedaj lahko uporabita kot skupni ključ za nadaljnjo komunikacijo.
\end{enumerate}

\vspace{0.5cm}

Oba sta izračunala enak rezultat, ker velja:
$$
A^b \equiv g^{ab} \equiv g^{ba} \equiv B^a \mod p
$$

\end{mdframed}

\newacronym{RSA}{RSA}{RSA}


\subsubsection{RSA}

Medtem ko prej opisana metoda omogoča samo dogovor o skupnem ključu, pa je bil RSA prvi pravi kriptosistem, ki je uporabljal javne ključe. Varnost RSA temelji na tem, da je faktorizacija produkta dveh velikih praštevil težko rešljiv problem za današnje računalnike.

Javni ključ je par $(n, e)$, zasebni ključ pa par $(n, d)$, ki ju dobimo po naslednjem postopku:

\begin{enumerate}
\item Izberemo dve naključni, približno enako veliki praštevili $p$ in $q$
\item Izračunamo $n=pq$
\item Izračunamo $\varphi(n)=\varphi(p)\varphi(q)=(p-1)(q-1)$
\item Izberemo število $e$, tako da $1<e<\varphi(n)$ tako da je $e$ tuje $\varphi(n)$
\item Izračunamo $d$, tako da $d \equiv e^{-1} \pmod{\varphi(n)}$
\end{enumerate}

Poleg zasebnega ključa moramo moramo paziti tudi na zasebnost vrednosti $p$, $q$ in $\varphi(n)$

\newacronym{ABE}{ABE}{Attribute Based Encryption}


\subsection{Šifriranje na osnovi pogojev}

Poglejmo si za primer Anito, ki shrani šifrirane podatke v oblak. Anita ima datoteko $X$, ki bi jo rada delila s skupinama $Prijatelji$ in $Znanci$. Prav tako ima datoteko $Y$, ki bi jo rada delila s skupinama $Družina$ in $Prijatelji$. To lahko stori tako, da ustvari en ključ za datoteko $X$ in ta ključ deli s $Prijatelji$ in $Znanci$, nato ustvari drug ključ za datoteko $Y$ in ta ključ deli s skupinama $Družina$ in $Prijatelji$. Problem je, da mora Anita za vsako datoteko, ki jo želi deliti, ustvariti nov ključ.

Naivna rešitev tega problema bi bila, da Anita ustvari en ključ za vsako skupino ljudi, s katerimi želi deliti datoteke. Vendar življenje ni preprosto, in lahko se zgodi, da neka oseba pripada večim skupinam s katerimi želi Anita deliti datoteke. Prav tako si lahko zaželi deliti neko datoteko samo z ljudmi, ki so hkrati $Prijatelji$ in $Sodelavci$. Recimo da je to datoteka $Z$. Tako datoteko lahko najprej šifrira s ključem, ki pripada skupini $Prijatelji$ in nato še s ključem, ki pripada skupini $Sodelavci$. Naivno bi mislili da je to dobra rešitev, vendar se izkaže da tak sistem ni odporen proti zaroti. Šifrirano datoteko lahko kot posamezniki odprejo samo ljudje, ki pripadajo obema skupinama, saj imajo le oni v posesti oba ključa. Vendar pa se lahko proti Aniti zarotita dva posameznika skupaj, kjer noben od nju ne pripada obema skupinama hkrati, vendar imata skupaj dostop do obeh ključev, saj vsak od nju pripada eni od skupin.

Ta problem nam rešuje šifriranje na podlagi pogojev, v nadaljevanju ABE. Kratica ABE pride iz angleške besede \gls{ABE}. ABE je primer asimetrične šifre, kjer je možno tajnopis dešifrirati z več različnimi ključi, pod pogojem da ima ključ take lastnosti oziroma atribute, kot so bili predpisani ob šifriranju.

Možnost dešifriranja določajo pravila, ki so predstavljena z drevesi, kjer so notranja vozlišča logične operacije, listi pa pogoji, ki so lahko zadovoljeni ali nezadovoljeni. Drevo beremo od listov proti korenu in za vsako notranje vozlišče ugotovimo ali je v zadovoljenem ali nezadovoljenem stanju. Na koncu pridemo do korena drevesa, in dešifriranje uspe samo v primeru, da je koren drevesa v zadovoljenem stanju.

Poznamo dve vrsti šifriranja na osnovi pogojev. V prvem je drevo, ki določa pravila shranjeno skupaj s ključem, v drugem primeru pa je shranjeno skupaj s tajnopisom.

Neki narobenapisan


\chapter{Uporaba kriptografije v oblaku}

Uvod, encryption at rest, in transit

\subsection{SSL in TLS}

\subsection{Server-side encryption}

\subsection{Client-side encryption}

\subsection{Homomorfne šifre}

\subsection{Attribute based encryption}

\subsection{Avtentikacija}




\newpage

%********************************************

\addcontentsline{toc}{chapter}{Seznam slik}
\addtocontents{toc}{\protect\vspace{-2ex}}
\listoffigures

\newpage

\addcontentsline{toc}{chapter}{Seznam tabel}
\listoftables

%\listofalgorithms


%********************************************

\newpage

\bibliographystyle{slplainurl}
\addcontentsline{toc}{chapter}{Literatura}
\label{stran_literatura}
\bibliography{diploma}


\end{document}
